\documentclass[../../main.tex]{subfiles}
\graphicspath{{../../image/}{../../figure/}}
\begin{document}

\section{粘性波扇}\label{section:ch09-粘性波扇}

\cref{section:ch08-黏性激波剖面}引入的粘性激波容许准则,把双曲守恒律方程组\eqref{align:ch09-9-1-1}的容许激波刻画为相伴耗散方程组\eqref{align:ch08-黏性方程}行波解当 $\mu\downarrow0$ 时的极限。此处的目标是把这一原理从单个激波推广到一般波扇。困难在于:与\eqref{align:ch09-9-1-1}不同,方程组\eqref{align:ch08-黏性方程}在时空坐标均匀拉伸下不变,因而不具有行波波扇解。作为补救,有人提出以带时变粘性的方程组代替\eqref{align:ch08-黏性方程},
\begin{align}
\partial_{t}U(x,t)+\partial_{x}F(U(x,t))=\mu t\,\partial_{x}^{2}U(x,t),
\label{align:ch09-9-8-1}
\end{align}
它在变换 $(x,t)\mapsto(\alpha x,\alpha t)$ 下不变。易见 $U=V_{\mu}(x/t)$ 是\eqref{align:ch09-9-8-1}的自相似解,当且仅当 $V_{\mu}(\xi)$ 满足常微分方程
\begin{align}
\mu\ddot{V}_{\mu}(\xi)=\dot{F}(V_{\mu}(\xi))-\xi\dot{V}_{\mu}(\xi).
\label{align:ch09-9-8-2}
\end{align}

\eqref{align:ch09-9-1-1}的自相似解 $U=V(x/t)$ 称为满足\emph{粘性波扇容许性准则},若 $V$ 是\eqref{align:ch09-9-8-2}的一族一致有界解 $V_{\mu}$ 当 $\mu\downarrow0$ 时的几乎处处极限。

除用作容许性检验外,粘性波扇准则还提示了构造 Riemann 问题\eqref{align:ch09-9-1-1}、\eqref{align:ch09-9-1-18}解的另一条途径。为此,须证明对任一固定 $\mu>0$,存在\eqref{align:ch09-9-8-2}在 $(-\infty,\infty)$ 上的某解 $V_{\mu}(\xi)$,附边值
\begin{align}
V_{\mu}(-\infty)=U_{L},\qquad V_{\mu}(+\infty)=U_{R},
\label{align:ch09-9-8-3}
\end{align}
然后证明族 $\{V_{\mu}(\xi):0<\mu<1\}$ 在 $(-\infty,\infty)$ 上有一致有界变差。此时由 Helly 定理(见\cref{section:ch01-BV函数理论概述}),可抽出收敛序列 $\{V_{\mu_{m}}\}$, $\mu_{m}\downarrow0$ ($m\to\infty$),其极限 $V$ 诱导出 Riemann 问题的解 $U(x,t)=V(x/t)$。

上述方案已在多种条件下成功实施。一种办法是对相当一般的数据 $U_{L}$、$U_{R}$ 求解 Riemann 问题,但限于特殊方程组,最著名的是两个守恒律的方程组。另一种办法是处理一般方程组,但限于中等强度波的情形,要求 $|U_{R}-U_{L}|$ 充分小。我们先考虑后一种情形。分析冗长而技术性强,故仅概述主要思想;细节读者可参阅\cref{section:ch09-注记}所引文献。

关键步骤是建立 $V_{\mu}(\xi)$ 在 $(-\infty,\infty)$ 上、不依赖于 $\mu$ 的总变差先验界。为给一般方程组做准备,先从标量守恒律\eqref{align:ch07-标量守恒律}入手。取 $\lambda(u)=f'(u)$ 并令 $\dot{V}_{\mu}(\xi)=a(\xi)$,把\eqref{align:ch09-9-8-2}写成
\begin{align}
\mu\dot{a}+[\xi-\lambda(V_{\mu}(\xi))]a=0.
\label{align:ch09-9-8-4}
\end{align}
\eqref{align:ch09-9-8-4}的解为 $a(\xi)=\tau\phi(\xi)$,其中
\begin{align}
\phi(\xi)=\frac{\exp[-\frac{1}{\mu}g(\xi)]}{\int_{-\infty}^{\infty}\exp[-\frac{1}{\mu}g(\zeta)]\,\mathrm{d}\zeta},
\label{align:ch09-9-8-5}
\end{align}
\begin{align}
g(\xi)=\int_{s}^{\xi}[\zeta-\lambda(V_{\mu}(\zeta))]\,\mathrm{d}\zeta.
\label{align:ch09-9-8-6}
\end{align}
积分下限 $s$ 取成使对一切 $\xi\in(-\infty,\infty)$ 有 $g(\xi)\geqslant0$。振幅 $\tau$ 由所赋边值确定,即 $V_{\mu}(-\infty)=u_{L}$, $V_{\mu}(\infty)=u_{R}$, $\tau=u_{R}-u_{L}$。由\eqref{align:ch09-9-8-5}知 $a(\xi)$ 的 $L^{1}$ 范数对 $\mu$ 一致有界,故族 $\{V_{\mu}:0<\mu<1\}$ 在 $(-\infty,\infty)$ 上一致有界变差。

现转到一般严格双曲方程组\eqref{align:ch09-9-1-1},把 $V_{\mu}(\xi)$ 实现为与不同特征族相联系的(复合)波的组装,写为
\begin{align}
\dot{V}_{\mu}(\xi)=\sum_{j=1}^{n}a_{j}(\xi)R_{j}(V_{\mu}(\xi)).
\label{align:ch09-9-8-7}
\end{align}
把\eqref{align:ch09-9-8-7}中的 $\dot{V}_{\mu}$ 代入\eqref{align:ch09-9-8-2},再从左端乘以 $L_{i}(V_{\mu}(\xi))$,推得
\begin{align}
\mu\dot{a}_{i}+[\xi-\lambda_{i}(V_{\mu}(\xi))]a_{i}=\mu\sum_{j,k=1}^{n}\beta_{ijk}(V_{\mu}(\xi))a_{j}a_{k},
\label{align:ch09-9-8-8}
\end{align}
其中
\begin{align}
\beta_{ijk}(U)=-L_{i}(U)\mathrm{D}R_{j}(U)R_{k}(U).
\label{align:ch09-9-8-9}
\end{align}
\eqref{align:ch09-9-8-8}左端与标量守恒律的\eqref{align:ch09-9-8-4}左端相同,右端则刻画不同特征族的相互作用。读者应注意\eqref{align:ch09-9-8-8}与\eqref{align:ch07-riccati}的类似性。还应注意,当方程组配有 Riemann 不变量坐标系统 $(w_{1},\dots,w_{n})$ 时, $a_{i}(\xi)=\dot{w}_{i}(V_{\mu}(\xi))$。此时如\cref{section:ch07-Riemann不变量}所示,当 $j\neq k$ 时 $\mathrm{D}R_{j}R_{k}$ 落在 $\{R_{j},R_{k}\}$ 张成的空间内,故\eqref{align:ch09-9-8-9}推出当 $i\neq j\neq k\neq i$ 时 $\beta_{ijk}=0$。对激波与稀疏波曲线重合的特殊方程组,如\eqref{align:ch07-电泳}, $\mathrm{D}R_{j}R_{j}$ 与 $R_{j}$ 共线,故即使 $i\neq j=k$ 也有 $\beta_{ijk}=0$,于是\eqref{align:ch09-9-8-8}中方程解耦。一般而言,分析的要旨是证明:在小振荡解的情形,即 $a_{i}$ 小,二阶相互作用效应更小。

\eqref{align:ch09-9-8-8}的解可分解为
\begin{align}
a_{i}(\xi)=\tau_{i}\phi_{i}(\xi)+\theta_{i}(\xi),
\label{align:ch09-9-8-10}
\end{align}
其中
\begin{align}
\phi_{i}(\xi)=\frac{\exp[-\frac{1}{\mu}g_{i}(\xi)]}{\int_{-\infty}^{\infty}\exp[-\frac{1}{\mu}g_{i}(\zeta)]\,\mathrm{d}\zeta},
\label{align:ch09-9-8-11}
\end{align}
\begin{align}
g_{i}(\xi)=\int_{s_{i}}^{\xi}[\zeta-\lambda_{i}(V_{\mu}(\zeta))]\,\mathrm{d}\zeta,
\label{align:ch09-9-8-12}
\end{align}
而 $\theta_{i}(\xi)$ 满足方程
\begin{align}
\mu\dot{\theta}_{i}+[\xi-\lambda_{i}(V_{\mu}(\xi))]\theta_{i}
=\mu\sum_{j,k=1}^{n}\beta_{ijk}(V_{\mu}(\xi))[\tau_{j}\phi_{j}(\xi)+\theta_{j}][\tau_{k}\phi_{k}(\xi)+\theta_{k}].
\label{align:ch09-9-8-13}
\end{align}

利用常数变易公式,微分方程\eqref{align:ch09-9-8-13}可等价改写为积分方程组:
\begin{align}
\theta_{i}(\xi)=\phi_{i}(\xi)\int_{c_{i}}^{\xi}\phi_{i}^{-1}(\zeta)\beta_{ijk}(V_{\mu}(\zeta))[\tau_{j}\phi_{j}(\zeta)+\theta_{j}(\zeta)][\tau_{k}\phi_{k}(\zeta)+\theta_{k}(\zeta)]\,\mathrm{d}\zeta.
\label{align:ch09-9-8-14}
\end{align}
仔细估计给出
\begin{align}
|\theta_{i}(\xi)|\leqslant c(\tau_{1}^{2}+\cdots+\tau_{n}^{2})\sum_{j=1}^{n}\phi_{j}(\xi),
\label{align:ch09-9-8-15}
\end{align}
这验证了在\eqref{align:ch09-9-8-10}中 $\theta_{i}$ 从属于 $\tau_{i}\phi_{i}$,即特征族在领头阶解耦。

可用压缩论证证明:对原点小邻域内任一固定 $(\tau_{1},\dots,\tau_{n})$,存在\eqref{align:ch09-9-8-2}在 $(-\infty,\infty)$ 上的某解 $V_{\mu}(\xi)$,满足\eqref{align:ch09-9-8-7}、\eqref{align:ch09-9-8-10}与\eqref{align:ch09-9-8-15}。为求解边值问题\eqref{align:ch09-9-8-2}、\eqref{align:ch09-9-8-3},须选取 $(\tau_{1},\dots,\tau_{n})$ 使
\begin{align}
\sum_{j=1}^{n}\int_{-\infty}^{\infty}[\tau_{j}\phi_{j}(\xi)+\theta_{j}(\xi)]R_{j}(V_{\mu}(\xi))\,\mathrm{d}\xi=U_{R}-U_{L}.
\label{align:ch09-9-8-16}
\end{align}
已证明\eqref{align:ch09-9-8-16}有唯一解 $(\tau_{1},\dots,\tau_{n})$,至少当 $|U_{R}-U_{L}|$ 充分小时。结果综述如下:

\begin{theorem}\label{theorem:ch09-粘性波扇存在性}
假定方程组\eqref{align:ch09-9-1-1}在 $\mathscr{O}$ 上严格双曲,固定任一状态 $U_{L}\in\mathscr{O}$。存在 $\delta>0$,使对任一满足 $|U_{R}-U_{L}|<\delta$ 的 $U_{R}\in\mathscr{O}$ 与每个 $\mu>0$,边值问题\eqref{align:ch09-9-8-2}、\eqref{align:ch09-9-8-3}都有解 $V_{\mu}(\xi)$,它有形如\eqref{align:ch09-9-8-7}、\eqref{align:ch09-9-8-10}的表示,其中 $(\tau_{1},\dots,\tau_{n})$ 接近原点,且 $\theta_{i}$ 满足\eqref{align:ch09-9-8-15}。此外,解族 $\{V_{\mu}(\xi):0<\mu<1\}$ 在 $(-\infty,\infty)$ 上有一致有界(且小)的全变差。特别地,可抽出序列 $\{V_{\mu_{m}}(\xi)\}$, $\mu_{m}\downarrow0$ ($m\to\infty$),它几乎处处有界收敛到函数 $V(\xi)$,使波扇 $U=V(x/t)$ 求解 Riemann 问题\eqref{align:ch09-9-1-1}、\eqref{align:ch09-9-1-18}。
\end{theorem}

仔细分析生成 $V(\xi)$ 作为序列 $\{V_{\mu_{m}}(\xi)\}$ 极限的过程,揭示 $V(\xi)$ 具有\cref{theorem:ch09-波扇结构}所述结构。此外,对 $V$ 的任一跳跃间断点 $\xi$,左侧的 $V(\xi-)$ 与右侧的 $V(\xi+)$ 由粘性激波剖面连接,因而满足粘性激波容许准则(取 $B=I$),如\cref{section:ch08-黏性激波剖面}所讨论。特别地, $V$ 的任一激波满足 Liu $E$ 条件,因而 $V$ 与\cref{theorem:ch09-严格双曲riemann解}所建立的唯一解重合。

\cref{theorem:ch09-粘性波扇存在性}所断言的边值问题\eqref{align:ch09-9-8-2}、\eqref{align:ch09-9-8-3}解 $V_{\mu}(\xi)$ 的构造以及族 $\{V_{\mu}\}$ 全变差界的推导,不依赖于方程组\eqref{align:ch09-9-8-1}是守恒的这一事实,而同样适用于任一方程组
\begin{align}
\mu\ddot{V}_{\mu}(\xi)=A(V_{\mu}(\xi))\dot{V}_{\mu}(\xi)-\xi\dot{V}_{\mu}(\xi),
\label{align:ch09-9-8-17}
\end{align}
只要矩阵 $A(U)$ 有实的互异特征值。若 $V(\xi)$ 是 $V_{\mu}(\xi)$ 的 $\mu\downarrow0$ 极限,则函数 $U(x,t)=V(x/t)$ 可解释为严格双曲、非守恒方程组
\begin{align}
\partial_{t}U+A(U)\partial_{x}U=0
\label{align:ch09-9-8-18}
\end{align}
的 Riemann 问题解,尽管它在分布意义下未必满足该方程组。

粘性波扇为一般严格双曲方程组\eqref{align:ch09-9-1-1}诱导出波曲线的另一种隐式构造,不要求分段真正非线性。

为追踪从某固定状态 $\bar{U}$ 发出的前向 $i$-波曲线,设右侧状态 $\hat{U}$ 与左侧 $\bar{U}$ 由一个中等强度 $i$-波连接。设此波是一族粘性波扇 $V_{\mu}(\xi)$ 的 $\mu\downarrow0$ 极限。于是 $V_{\mu}$ 定义在 $\bar{\xi}=\lambda_{i}(\bar{U})$ 的小邻域内,取值接近 $\bar{U}$,且 $\mu\dot{V}_{\mu}(\xi)$ 小。我们通过重新标度变量 $\xi=\mu\zeta$ 拉伸定义域。再把 $a_{j}$ 在展开\eqref{align:ch09-9-8-7}中重新标度,令 $w_{j}=\mu a_{j}$,并组装向量 $W=(w_{1},\dots,w_{n})$。于是可把\eqref{align:ch09-9-8-7}、\eqref{align:ch09-9-8-8}改写为自治一阶方程组
\begin{align}
\begin{cases}
V'=\displaystyle\sum_{j=1}^{n}w_{j}R_{j}(V),\\[2mm]
w_{j}'=[\lambda_{j}(V)-\xi]w_{j}+\displaystyle\sum_{k,\ell=1}^{n}\beta_{jk\ell}(V)w_{k}w_{\ell},&j=1,\dots,n,\\
\xi'=\mu,\\
\mu'=0,
\end{cases}
\label{align:ch09-9-8-19}
\end{align}
其中撇表示对 $\zeta$ 求导。

在平衡点 $V=\bar{U}$、$W=0$、$\xi=\lambda_{i}(\bar{U})$、$\mu=0$ 附近线性化\eqref{align:ch09-9-8-19},得方程组
\begin{align}
\begin{cases}
V'=\displaystyle\sum_{j=1}^{n}w_{j}R_{j}(\bar{U}),\\[2mm]
w_{j}'=[\lambda_{j}(\bar{U})-\xi]w_{j},&j=1,\dots,n,\\
\xi'=\mu,\\
\mu'=0.
\end{cases}
\label{align:ch09-9-8-20}
\end{align}

该系统的中心子空间 $\mathscr{N}$ 由所有满足 $j\neq i$ 时 $w_{j}=0$ 的向量 $(V,W,\xi,\mu)\in\mathbb{R}^{2n+2}$ 组成,维数为 $n+3$。由中心流形定理,停留于上述平衡点邻域内的\eqref{align:ch09-9-8-19}任一解必落在一个 $(n+3)$ 维流形 $\mathscr{M}$ 上,该流形在平衡点处与 $\mathscr{N}$ 相切,在\eqref{align:ch09-9-8-19}生成的流下不变,并有局部表示
\begin{align}
w_{j}=\varphi_{j}(V,\omega,\xi;\mu),\qquad j\neq i,
\label{align:ch09-9-8-21}
\end{align}
其中 $\omega$ 表示 $w_{i}$。由斜积流理论,可选取函数 $\varphi_{j}$ 使对一切接近 $\bar{U}$、$\lambda_{i}(\bar{U})$、$0$ 的 $V$、$\xi$、$\mu$,有 $\varphi_{j}(V,0,\xi;\mu)=0$。于是可置
\begin{align}
w_{j}=\omega\psi_{j}(V,\omega,\xi;\mu),\qquad j\neq i,
\label{align:ch09-9-8-22}
\end{align}
其中 $\psi_{j}(\bar{U},0,\lambda_{i}(\bar{u});0)=0$,因为 $\mathscr{M}$ 在平衡点处与 $\mathscr{N}$ 相切。再引入新变量 $\tau$ 使
\begin{align}
\frac{\mathrm{d}}{\mathrm{d}\tau}=\frac{1}{\omega}\frac{\mathrm{d}}{\mathrm{d}\zeta}.
\label{align:ch09-9-8-23}
\end{align}

为看清解的分量 $(V_{\mu},\omega_{\mu},\xi_{\mu})$ 作为 $\tau$ 的函数在 $\mathscr{M}$ 上如何演化,合并\eqref{align:ch09-9-8-19}、\eqref{align:ch09-9-8-22}与\eqref{align:ch09-9-8-23},推得
\begin{align}
\frac{\mathrm{d}V_{\mu}}{\mathrm{d}\tau}=P_{\mu}(V_{\mu},\omega_{\mu},\xi_{\mu}),
\label{align:ch09-9-8-24}
\end{align}
\begin{align}
\frac{\mathrm{d}\omega_{\mu}}{\mathrm{d}\tau}=p_{\mu}(V_{\mu},\omega_{\mu},\xi_{\mu})-\xi_{\mu},
\label{align:ch09-9-8-25}
\end{align}
其中已置
\begin{align}
P_{\mu}(V,\omega,\xi)=R_{i}(V)+\sum_{j\neq i}\psi_{j}(V,\omega,\xi;\mu)R_{j}(V),
\label{align:ch09-9-8-26}
\end{align}
\begin{align}
p_{\mu}(V,\omega,\xi)=\lambda_{i}(V)+\sum_{k,\ell=1}^{n}\omega\beta_{ik\ell}(V)\psi_{k}(V,\omega,\xi;\mu)\psi_{\ell}(V,\omega,\xi;\mu).
\label{align:ch09-9-8-27}
\end{align}
特别地, $P_{0}(\bar{U},0,\lambda_{i}(\bar{U}))=R_{i}(\bar{U})$, $p_{0}(\bar{U},0,\lambda_{i}(\bar{U}))=\lambda_{i}(\bar{U})$。

为推导 $\xi_{\mu}(\tau)$ 的方程,注意\eqref{align:ch09-9-8-23}与\eqref{align:ch09-9-8-19}$_{3}$ 合起来给出 $\mathrm{d}\xi_{\mu}/\mathrm{d}\tau=\xi_{\mu}'/\omega=\mu/\omega$。我们对该关系关于 $\tau$ 求导,并用\eqref{align:ch09-9-8-25}得
\begin{align}
\mu\frac{\mathrm{d}^{2}\xi_{\mu}}{\mathrm{d}\tau^{2}}=-\left(\frac{\mathrm{d}\xi_{\mu}}{\mathrm{d}\tau}\right)^{2}[p_{\mu}(V_{\mu},\omega_{\mu},\xi_{\mu})-\xi_{\mu}].
\label{align:ch09-9-8-28}
\end{align}

当 $\mu\downarrow0$ 时, $(V_{\mu},\omega_{\mu},\xi_{\mu})$ 一致收敛到 $(V,\omega,\xi)$。特别地, $V(0)=\bar{U}$, $\omega(0)=0$,且 $V(s)=\hat{U}$,其中 $s$ 不妨取为正、小。由\eqref{align:ch09-9-8-28}, $[0,s]$ 是与激波相联系、其上 $\mathrm{d}\xi/\mathrm{d}\tau=0$ 的 $\tau$-区间之至多可数并,与及稀疏波相联系、其上 $\xi=p_{0}(V,\omega,\xi)$ 的 $\tau$-区间之并。此外,在从激波到稀疏波(或从稀疏波到激波)的转换点处, $\mathrm{d}^{2}\xi_{\mu}/\mathrm{d}\tau^{2}$ 应非负(或非正)。于是推得
\begin{align}
\xi(\tau)=\frac{\mathrm{d}g}{\mathrm{d}\tau}(\tau),\qquad 0\leqslant\tau\leqslant s,
\label{align:ch09-9-8-29}
\end{align}
其中 $g$ 是函数
\begin{align}
f(\tau)=\int_{0}^{\tau}p_{0}(V(\sigma),\omega(\sigma),\xi(\sigma))\,\mathrm{d}\sigma,\qquad 0\leqslant\tau\leqslant s,
\label{align:ch09-9-8-30}
\end{align}
在 $[0,s]$ 上的凸包,即 $g(\tau)=\inf\{\theta_{1}f(\tau_{1})+\theta_{2}f(\tau_{2}):\theta_{1}\geqslant0,\theta_{2}\geqslant0,\theta_{1}+\theta_{2}=1,\ 0\leqslant\tau_{1}\leqslant\tau_{2}\leqslant s,\ \theta_{1}\tau_{1}+\theta_{2}\tau_{2}=\tau\}$。然后\eqref{align:ch09-9-8-24}与\eqref{align:ch09-9-8-25}给出
\begin{align}
V(\tau)=\bar{U}+\int_{0}^{\tau}P_{0}(V(\sigma),\omega(\sigma),\xi(\sigma))\,\mathrm{d}\sigma,\qquad 0\leqslant\tau\leqslant s,
\label{align:ch09-9-8-31}
\end{align}
\begin{align}
\omega(\tau)=f(\tau)-g(\tau),\qquad 0\leqslant\tau\leqslant s.
\label{align:ch09-9-8-32}
\end{align}

可以证明,一旦 $P_{0}(V,\omega,\xi)$ 与 $p_{0}(V,\omega,\xi)$ 指定,方程组\eqref{align:ch09-9-8-29}、\eqref{align:ch09-9-8-31}、\eqref{align:ch09-9-8-32}可用 Picard 迭代求解,对任意小正数 $s$ 给出 $[0,s]$ 上的函数 $V(\tau)$、$\omega(\tau)$、$\xi(\tau)$。负 $s$ 的处理类似,只是此时 $g$ 是 $f$ 在 $[s,0]$ 上的凹包。因而这些方程提供了从 $\bar{U}$ 发出的 $i$-波曲线 $\Phi_{i}$ 的一个隐式表示,只需置 $\Phi_{i}(s;\bar{U})=V(s)$。由\eqref{align:ch09-9-8-23}的定义, $\tau$ 近似等于 $V-\bar{U}$ 在 $R_{i}$ 上的投影。因而上述 $i$-波曲线的构造与\cref{section:ch09-一般严格双曲方程组}开头描述的标量守恒律波构造极为相似。

接下来的课题是用粘性波扇法构造两个守恒律方程组 Riemann 问题的解,
\begin{align}
\begin{cases}
\partial_{t}u+\partial_{x}f(u,v)=0,\\
\partial_{t}v+\partial_{x}g(u,v)=0,
\end{cases}
\label{align:ch09-9-8-33}
\end{align}
但初值不受限制
\begin{align}
(u(x,0),v(x,0))=
\begin{cases}
(u_{L},v_{L}),&x<0,\\
(u_{R},v_{R}),&x>0.
\end{cases}
\label{align:ch09-9-8-34}
\end{align}
关键限制是 $f_{v}$ 与 $g_{u}$ 同号,不妨设
\begin{align}
f_{v}(u,v)<0,\qquad g_{u}(u,v)<0,\qquad\text{对一切 }(u,v).
\label{align:ch09-9-8-35}
\end{align}
特别地,该方程组严格双曲。对称耦合方程组与等熵弹性动力学方程组\eqref{align:ch07-p系统}是此类型的典型代表。分析将展示 $\delta$ 激波如何作为粘性剖面极限中的"集中"而出现。

此处\eqref{align:ch09-9-8-2}、\eqref{align:ch09-9-8-3}取形式
\begin{align}
\begin{cases}
\mu\ddot{u}_{\mu}(\xi)=\dot{f}(u_{\mu}(\xi),v_{\mu}(\xi))-\xi\dot{u}_{\mu}(\xi),\\
\mu\ddot{v}_{\mu}(\xi)=\dot{g}(u_{\mu}(\xi),v_{\mu}(\xi))-\xi\dot{v}_{\mu}(\xi),
\end{cases}
\label{align:ch09-9-8-36}
\end{align}
\begin{align}
(u_{\mu}(-\infty),v_{\mu}(-\infty))=(u_{L},v_{L}),\qquad(u_{\mu}(\infty),v_{\mu}(\infty))=(u_{R},v_{R}).
\label{align:ch09-9-8-37}
\end{align}

假设\eqref{align:ch09-9-8-35}的重要性由下面引理可见。

\begin{lemma}\label{lemma:ch09-粘性解单调性}
设 $(u_{\mu}(\xi),v_{\mu}(\xi))$ 是\eqref{align:ch09-9-8-36}与\eqref{align:ch09-9-8-37}在 $(-\infty,\infty)$ 上的解。则下列之一成立:
\begin{enumerate}
\item $u_{\mu}(\xi)$ 与 $v_{\mu}(\xi)$ 在 $(-\infty,\infty)$ 上均为常值。
\item $u_{\mu}(\xi)$ 严格递增(或递减),在 $(-\infty,\infty)$ 上无临界点; $v_{\mu}(\xi)$ 在 $(-\infty,\infty)$ 上至多有一个临界点,且必为极大(或极小)。
\item $v_{\mu}(\xi)$ 严格递增(或递减),在 $(-\infty,\infty)$ 上无临界点; $u_{\mu}(\xi)$ 在 $(-\infty,\infty)$ 上至多有一个临界点,且必为极大(或极小)。
\end{enumerate}
\end{lemma}

\begin{proof}
注意 $\dot{u}_{\mu}(\xi_{0})=0$ 与 $\ddot{u}_{\mu}(\xi_{0})=0$ 推出 $\dot{v}_{\mu}(\xi_{0})=0$;而 $\dot{v}_{\mu}(\xi_{0})=0$ 与 $\ddot{v}_{\mu}(\xi_{0})=0$ 推出 $\dot{u}_{\mu}(\xi_{0})=0$。因而,由常微分方程初值问题解的唯一性,若 $u_{\mu}(\xi)$ 与 $v_{\mu}(\xi)$ 中之一有退化临界点,则两函数在 $(-\infty,\infty)$ 上必都为常值。

转到非退化临界点,注意 $\dot{u}_{\mu}(\xi_{0})=0$ 与 $\ddot{u}_{\mu}(\xi_{0})<0$(或 $>0$)推出 $\dot{v}_{\mu}(\xi_{0})>0$(或 $<0$);类似地, $\dot{v}_{\mu}(\xi_{0})=0$ 与 $\ddot{v}_{\mu}(\xi_{0})<0$(或 $>0$)推出 $\dot{u}_{\mu}(\xi_{0})>0$(或 $\dot{u}_{\mu}(\xi_{0})<0$)。

现设 $v_{\mu}(\xi)$ 有不止一个非退化临界点,取相邻两个, $\xi_{1}$ 处为极大、$\xi_{2}$ 处为极小。不妨设 $\xi_{1}<\xi_{2}$。则 $\dot{v}_{\mu}(\xi_{1})=0$, $\ddot{v}_{\mu}(\xi_{1})<0$, $\dot{v}_{\mu}(\xi_{2})=0$, $\ddot{v}_{\mu}(\xi_{2})>0$,且对 $\xi\in(\xi_{1},\xi_{2})$ 有 $\dot{v}_{\mu}(\xi)<0$。于是 $\dot{u}_{\mu}(\xi_{1})>0$, $\dot{u}_{\mu}(\xi_{2})<0$。故存在 $\xi_{0}\in(\xi_{1},\xi_{2})$ 使 $\dot{u}_{\mu}(\xi_{0})=0$ 且 $\ddot{u}_{\mu}(\xi_{0})<0$。但这推出 $\dot{v}_{\mu}(\xi_{0})>0$,矛盾。$\xi_{1}>\xi_{2}$ 的情形同样矛盾。同一论证表明 $u_{\mu}(\xi)$ 至多有一个非退化临界点。

最后,设 $u_{\mu}(\xi)$ 与 $v_{\mu}(\xi)$ 都有非退化临界点,分别在 $\xi_{1}$ 与 $\xi_{2}$。不妨设 $\xi_{1}<\xi_{2}$ 且 $\xi_{2}$ 是 $v_{\mu}(\xi)$ 的极大点。则对 $\xi\in(-\infty,\xi_{2})$ 有 $\dot{v}_{\mu}(\xi)>0$,而 $\dot{v}_{\mu}(\xi_{2})=0$, $\ddot{v}_{\mu}(\xi_{2})<0$。这推出 $\dot{u}_{\mu}(\xi_{2})>0$,从而 $\xi_{1}$ 必为 $u_{\mu}(\xi)$ 的极小点, $\dot{u}_{\mu}(\xi_{1})=0$, $\ddot{u}_{\mu}(\xi_{1})>0$。这又推出 $\dot{v}_{\mu}(\xi_{1})<0$,矛盾。其他可能组合导致类似矛盾。证明完毕。
\end{proof}

由于 $u_{\mu}(\xi)$ 与 $v_{\mu}(\xi)$ 图像的构型非常特殊,建立\eqref{align:ch09-9-8-36}、\eqref{align:ch09-9-8-37}解的存在性相对容易。事实上,结果表明为此只需用给定数据 $(u_{L},v_{L})$、$(u_{R},v_{R})$ 与参数 $\mu$,先验界定 $u_{\mu}(\xi)$ 或 $v_{\mu}(\xi)$ 所取的唯一"峰"值即可。这类估计的推导及由此得到的存在性证明,读者可在\cref{section:ch09-注记}所引文献中找到,其假设是:要么 $f(u,v)$、$g(u,v)$ 的增长受限于
\begin{align}
|f(u,v)|\leqslant h(v)(1+|u|)^{p},\qquad|g(u,v)|\leqslant h(u)(1+|v|)^{p},
\label{align:ch09-9-8-38}
\end{align}
其中 $h$ 是连续函数, $p<2$;要么方程组\eqref{align:ch09-9-8-33}配有一个熵 $\eta(u,v)$,其 Hessian 矩阵 $\mathrm{D}^{2}\eta(u,v)$ 的特征值下方受 $(1+|u|)^{-p}(1+|v|)^{-p}$ 控制,某 $p<3$。第一类特别包含\eqref{align:ch07-p系统},第二类包含所有对称方程组。

设 $(u_{\mu},v_{\mu})$ 存在,令 $\mu\downarrow0$ 取极限以得到 Riemann 问题的解。为此需不依赖于 $\mu$ 的估计。不妨设 $v_{\mu}(\xi)$ 在 $(-\infty,\infty)$ 上严格递增,而 $u_{\mu}(\xi)$ 在 $(-\infty,\xi_{\mu})$ 上严格递增,在 $\xi_{\mu}$ 处达最大,在 $(\xi_{\mu},\infty)$ 上严格递减。其他可能构型类似处理。

令 $\bar{u}=\max\{u_{L},u_{R}\}$,并取点 $\xi_{\ell}\in(-\infty,\xi_{\mu})\cup\{-\infty\}$ 与 $\xi_{r}\in(\xi_{\mu},\infty)\cup\{\infty\}$ 使 $u(\xi_{\ell})=u(\xi_{r})=\bar{u}$。对任一开区间 $(a,b)\subset(-\infty,\infty)$,利用\eqref{align:ch09-9-8-36}$_{1}$ 与\eqref{align:ch09-9-8-35},推得
\begin{align}
\int_{a}^{b}[u_{\mu}(\xi)-\bar{u}]\,\mathrm{d}\xi
\leqslant\int_{\xi_{\ell}}^{\xi_{r}}[u_{\mu}(\xi)-\bar{u}]\,\mathrm{d}\xi
=-\int_{\xi_{\ell}}^{\xi_{r}}\xi\dot{u}_{\mu}(\xi)\,\mathrm{d}\xi
=\mu\dot{u}_{\mu}(\xi_{r})-\mu\dot{u}_{\mu}(\xi_{\ell})-f(\bar{u},v_{\mu}(\xi_{r}))+f(\bar{u},v_{\mu}(\xi_{\ell}))
\leqslant f(\bar{u},v_{L})-f(\bar{u},v_{R}).
\label{align:ch09-9-8-39}
\end{align}

由\eqref{align:ch09-9-8-39},存在序列 $\{\mu_{k}\}$, $\mu_{k}\downarrow0$ ($k\to0$),使 $\{\xi_{\mu_{k}}\}$ 收敛到某点 $\xi_{0}\in(-\infty,\infty)\cup\{\pm\infty\}$, $\{v_{\mu_{k}}(\xi)\}$ 在 $(-\infty,\infty)$ 上点点收敛到某单调递增函数 $v(\xi)$,而 $\{u_{\mu_{k}}(\xi)\}$ 在 $(-\infty,\xi_{0})\cup(\xi_{0},\infty)$ 上点点收敛到某局部可积函数 $u(\xi)$,它在 $(-\infty,\xi_{0})$ 上单调递增、在 $(\xi_{0},\infty)$ 上单调递减。此外易见 $u(-\infty)=u_{L}$, $u(\infty)=u_{R}$, $v(-\infty)=v_{L}$, $v(\infty)=v_{R}$。

当 $\xi_{0}=-\infty$(或 $\xi_{0}=\infty$)时, $u(\xi)$ 在 $(-\infty,\infty)$ 上单调递增(或递减),此时 $(u(\xi),v(\xi))$ 是 Riemann 问题的标准解。当 $\xi_{0}\in(-\infty,\infty)$ 时情形变得有趣:此时当 $k\to\infty$ 时 $u_{\mu_{k}}\to u+\omega\delta_{\xi_{0}}$,在分布意义下,其中 $\delta_{\xi_{0}}$ 表示 $\xi_{0}$ 处的 Dirac $\delta$ 函数, $\omega\geqslant0$。

我们以检验函数 $\varphi\in C_{0}^{\infty}(-\infty,\infty)$ 乘\eqref{align:ch09-9-8-36}两式,在 $(-\infty,\infty)$ 上积分并分部积分,得
\begin{align}
\begin{cases}
\displaystyle\int_{-\infty}^{\infty}\{\mu u_{\mu}\ddot{\varphi}+[f(u_{\mu},v_{\mu})-\xi u_{\mu}]\dot{\varphi}-u_{\mu}\varphi\}\,\mathrm{d}\xi=0,\\[2mm]
\displaystyle\int_{-\infty}^{\infty}\{\mu v_{\mu}\ddot{\varphi}+[g(u_{\mu},v_{\mu})-\xi v_{\mu}]\dot{\varphi}-v_{\mu}\varphi\}\,\mathrm{d}\xi=0.
\end{cases}
\label{align:ch09-9-8-40}
\end{align}
对在含 $\xi_{0}$ 的某开区间上为常数的检验函数施用\eqref{align:ch09-9-8-40},并沿序列 $\{\mu_{k}\}$ 令 $\mu\downarrow0$,得
\begin{align}
\begin{cases}
\displaystyle\int_{-\infty}^{\infty}[f(u(\xi),v(\xi))-\xi u(\xi)]\dot{\varphi}(\xi)\,\mathrm{d}\xi
=\int_{-\infty}^{\infty}u(\xi)\varphi(\xi)\,\mathrm{d}\xi+\omega\varphi(\xi_{0}),\\[2mm]
\displaystyle\int_{-\infty}^{\infty}[g(u(\xi),v(\xi))-\xi v(\xi)]\dot{\varphi}(\xi)\,\mathrm{d}\xi
=\int_{-\infty}^{\infty}v(\xi)\varphi(\xi)\,\mathrm{d}\xi.
\end{cases}
\label{align:ch09-9-8-41}
\end{align}
把 $\varphi$ 的支集在 $\xi_{0}$ 周围收缩,推出
\begin{align}
\begin{cases}
\lim_{\xi\uparrow\xi_{0}}[f(u(\xi),v(\xi))-\xi u(\xi)]-\lim_{\xi\downarrow\xi_{0}}[f(u(\xi),v(\xi))-\xi u(\xi)]=\omega,\\[1mm]
\lim_{\xi\uparrow\xi_{0}}[g(u(\xi),v(\xi))-\xi v(\xi)]-\lim_{\xi\downarrow\xi_{0}}[g(u(\xi),v(\xi))-\xi v(\xi)]=0,
\end{cases}
\label{align:ch09-9-8-42}
\end{align}
其中四个极限都(有限)存在。特别地,这推出函数 $f(u(\xi),v(\xi))$ 与 $g(u(\xi),v(\xi))$ 在 $(-\infty,\infty)$ 上局部可积,且\eqref{align:ch09-9-8-41}对任意 $\varphi\in C_{0}^{\infty}(-\infty,\infty)$ 成立。等价地,
\begin{align}
\begin{cases}
[f(u,v)-\xi u]^{\cdot}+u+\omega\delta_{\xi_{0}}=0,\\
[g(u,v)-\xi v]^{\cdot}+v=0,
\end{cases}
\label{align:ch09-9-8-43}
\end{align}
在分布意义下。于是我们断定:当 $\omega=0$ 时 $(u(\xi),v(\xi))$ 是 Riemann 问题的标准解,可能有 $u(\xi_{0})=\infty$,正如方程组\eqref{align:ch07-p系统}的解\eqref{align:ch09-9-8-43};而当 $\omega>0$ 时, $(u(\xi)+\omega\delta_{\xi_{0}},v(\xi))$ 可解释为 Riemann 问题的非标准解,含一个位于 $\xi_{0}$ 的 $\delta$ 激波,正如方程组\eqref{align:ch07-p系统}的解\eqref{align:ch09-9-6-6}。

现假定方程组配有一个凸熵 $\eta(u,v)$ 的熵-熵流对 $(\eta,q)$,具超线性增长,
\begin{align}
\frac{\eta(u,v)}{|u|+|v|}\to\infty,\qquad\text{当 }|u|+|v|\to\infty,
\label{align:ch09-9-8-44}
\end{align}
并证明 $(u(\xi),v(\xi))$ 是 Riemann 问题的标准解,即 $\omega=0$。以 $\eta_{u}(u_{\mu},v_{\mu})$ 乘\eqref{align:ch09-9-8-36}$_{1}$,以 $\eta_{v}(u_{\mu},v_{\mu})$ 乘\eqref{align:ch09-9-8-36}$_{2}$,把两式相加得
\begin{align}
\mu\ddot{\eta}(u_{\mu},v_{\mu})-\mu[\eta_{uu}\dot{u}_{\mu}^{2}+2\eta_{uv}\dot{u}_{\mu}\dot{v}_{\mu}+\eta_{vv}\dot{v}_{\mu}^{2}]
=\dot{q}(u_{\mu},v_{\mu})-\xi\dot{\eta}(u_{\mu},v_{\mu}).
\label{align:ch09-9-8-45}
\end{align}
令 $\bar{\eta}=\max\{\eta(u_{L},v_{L}),\eta(u_{R},v_{R})\}$,并取 $(-\infty,\xi_{0})\cup\{-\infty\}$ 中最大数 $\xi_{L}$ 与 $(\xi_{0},\infty)\cup\{\infty\}$ 中最小数 $\xi_{R}$,使 $\eta(u_{\mu}(\xi_{L}),v_{\mu}(\xi_{L}))=\eta(u_{\mu}(\xi_{R}),v_{\mu}(\xi_{R}))=\bar{\eta}$。利用\eqref{align:ch09-9-8-45},
\begin{align}
\int_{\xi_{L}}^{\xi_{R}}[\eta(u_{\mu},v_{\mu})-\bar{\eta}]\,\mathrm{d}\xi
=-\int_{\xi_{L}}^{\xi_{R}}\xi\dot{\eta}(u_{\mu},v_{\mu})\,\mathrm{d}\xi
\leqslant q(u_{\mu}(\xi_{L}),v_{\mu}(\xi_{L}))-q(u_{\mu}(\xi_{R}),v_{\mu}(\xi_{R})).
\label{align:ch09-9-8-46}
\end{align}

\eqref{align:ch09-9-8-46}右端对 $\mu>0$ 一致有界。因而合并\eqref{align:ch09-9-8-46}与\eqref{align:ch09-9-8-44}得
\begin{align}
\int_{\{u_{\mu}\geqslant\bar{u}\}}u_{\mu}(\xi)\,\mathrm{d}\xi\to0,\qquad\text{当 }\bar{u}\to\infty,
\label{align:ch09-9-8-47}
\end{align}
对 $\mu>0$ 一致,从而 $\omega=0$。

显然同一论证适用于 $(u_{\mu}(\xi),v_{\mu}(\xi))$ 的一切可能构型。于是我们建立了

\begin{theorem}\label{theorem:ch09-两方程粘性极限}
假定方程组\eqref{align:ch09-9-8-33}(其中 $f_{v}g_{u}>0$)配有一个凸熵 $\eta(u,v)$,具超线性增长\eqref{align:ch09-9-8-44}。则\eqref{align:ch09-9-8-36}、\eqref{align:ch09-9-8-37}的解序列 $\{(u_{\mu_{k}},v_{\mu_{k}})\}$, $\mu_{k}\to0$ ($k\to\infty$),点点收敛,且在分布意义下收敛到 Riemann 问题\eqref{align:ch09-9-8-33}、\eqref{align:ch09-9-8-34}的标准解 $(u,v)$。$u(\xi)$、$v(\xi)$ 中至少一个在 $(-\infty,\infty)$ 上单调,另一个至多有一个极值,该极值可以有界或无界。
\end{theorem}

特别地,任一 $f_{v}=g_{u}\neq0$ 的两个守恒律对称方程组都满足上述定理的假设。在\cref{section:ch09-注记}所引文献中,读者将找到关于 $f$、$g$ 的假设,使所得 Riemann 问题解必有界。还已证明,这些解中任一激波满足粘性激波容许准则,从而满足 Liu $E$ 条件。

承接\cref{section:ch09-注记}的讨论,人们可以论证:端状态为 $U_{L}$、$U_{R}$ 的 Riemann 问题波扇解,除非能在抛物方程组\cref{section:ch09-注记}的解的 $t\to\infty$ 渐近中,以衰减到 $U_{L}$、$U_{R}$ 的初值 $U_{0}(x)$ 当 $x\to\mp\infty$ 时捕捉到,否则不应称为容许。事实上,\cref{section:ch09-注记}中关于粘性激波剖面渐近稳定性的结果处理了上述问题的一个特殊情形。其互补特殊情形,即稀疏波的渐近稳定性,也已被广泛研究(文献见\cref{section:ch09-注记})。把上述两个成分结合起来以合成 Riemann 问题完整解的任务,尚未以决定性方式完成。

\end{document}
